% !TEX root =./ECE343_Practice_main.tex
%
% L17 practice set (Vector and Coordinate System Review): differential
% elements and integrals, gradient, divergence, curl, the Laplacian, and the
% divergence and Stokes' theorems.

\fullwidth{\textbf{Learning Objectives 4.1--4.3: } I can apply the gradient,
divergence, curl, and Laplacian in the three major coordinate systems; calculate
surface and volume integrals in each; and apply the divergence theorem and
Stokes' theorem to simplify integrals.
\\[4pt]\normalfont\small Show your work. Write a \textbf{Documentation} line at the
end (\textbf{None} if you did not collaborate).}

\begin{questions}

%%%%%%%%%%%%%%%%%%%%%%%%%%%%%%%%%%%%%%%%%
\question \textbf{Differential elements and integrals.}
\begin{parts}

	\part Starting from the distance a point moves when each spherical coordinate
	changes by a small amount, derive the spherical volume element $dv$.
	\begin{solution}[1.3in]
	Changing $r$ by $dr$ moves the point $dr$. Changing $\theta$ by $d\theta$
	moves it along a circle of radius $r$, a distance $r\,d\theta$. Changing
	$\phi$ by $d\phi$ moves it along a horizontal circle of radius $r\sin\theta$,
	a distance $r\sin\theta\,d\phi$. The volume element is the product of the
	three edges:
	\eq{ dv = (dr)(r\,d\theta)(r\sin\theta\,d\phi) = r^2\sin\theta\,dr\,d\theta\,d\phi }
	\end{solution}

	\begin{minipage}[t]{0.58\linewidth}
		\part Find the area of the cap of a sphere of radius $a$ that lies within
		$0 \le \theta \le \theta_0$. Check your answer at $\theta_0 = \pi$.
		\begin{solution}[1.3in]
		\eq{
			S &= \int_0^{2\pi}\!\!\int_0^{\theta_0} a^2\sin\theta\,d\theta\,d\phi
			= 2\pi a^2\left[-\cos\theta\right]_0^{\theta_0} \\
			&= 2\pi a^2\left(1 - \cos\theta_0\right)
		}
		At $\theta_0 = \pi$, $S = 4\pi a^2$, the whole sphere.
		\end{solution}
	\end{minipage}\hfill%
	\begin{minipage}[t]{0.37\linewidth}
		\raggedleft \vspace{12pt}
		$S$ = \ansbox[1.5in]{$2\pi a^2(1 - \cos\theta_0)$}
	\end{minipage}

	\begin{minipage}[t]{0.58\linewidth}
		\part Find the volume of a spherical shell with inner radius $a$ and outer
		radius $b$.
		\begin{solution}[1.1in]
		\eq{
			V &= \int_0^{2\pi}\!\!\int_0^{\pi}\!\!\int_a^{b} r^2\sin\theta\,dr\,d\theta\,d\phi \\
			&= 2\pi \cdot 2 \cdot \dfrac{b^3 - a^3}{3} = \dfrac{4\pi}{3}\left(b^3 - a^3\right)
		}
		\end{solution}
	\end{minipage}\hfill%
	\begin{minipage}[t]{0.37\linewidth}
		\raggedleft \vspace{12pt}
		$V$ = \ansbox[1.5in]{$\tfrac{4\pi}{3}(b^3 - a^3)$}
	\end{minipage}

\end{parts}

\newpage
%%%%%%%%%%%%%%%%%%%%%%%%%%%%%%%%%%%%%%%%%
\question \textbf{Gradient, divergence, and the Laplacian.}
\begin{parts}

	\begin{minipage}[t]{0.58\linewidth}
		\part Find the gradient of $f = xy$, and its magnitude at the point
		$(-4, 5, 0)$.
		\begin{solution}[1in]
		\eq{
			\nabla f &= \hat{\mathbf{x}}\,y + \hat{\mathbf{y}}\,x
			= 5\hat{\mathbf{x}} - 4\hat{\mathbf{y}} \\
			|\nabla f| &= \sqrt{5^2 + 4^2} = \sqrt{41} = 6.40
		}
		\end{solution}
	\end{minipage}\hfill%
	\begin{minipage}[t]{0.37\linewidth}
		\raggedleft \vspace{12pt}
		$|\nabla f|$ = \ansbox[1.3in]{$6.40$}
	\end{minipage}

	\begin{minipage}[t]{0.58\linewidth}
		\part Find the divergence of
		$\mathbf{F} = \hat{\boldsymbol{\rho}}\,\rho^2 + \hat{\mathbf{z}}\,z$.
		\begin{solution}[1in]
		\eq{
			\nabla\cdot\mathbf{F} = \dfrac{1}{\rho}\dfrac{\partial\left(\rho\cdot\rho^2\right)}{\partial\rho} + \dfrac{\partial z}{\partial z}
			= \dfrac{3\rho^2}{\rho} + 1 = 3\rho + 1
		}
		\end{solution}
	\end{minipage}\hfill%
	\begin{minipage}[t]{0.37\linewidth}
		\raggedleft \vspace{12pt}
		$\nabla\cdot\mathbf{F}$ = \ansbox[1.3in]{$3\rho + 1$}
	\end{minipage}

	\begin{minipage}[t]{0.58\linewidth}
		\part Find the Laplacian of $f = e^{xyz}$.
		\begin{solution}[1in]
		\eq{
			\dfrac{\partial^2 f}{\partial x^2} &= y^2z^2 e^{xyz}, \quad \text{and likewise in } y \text{ and } z \\
			\nabla^2 f &= \left(y^2z^2 + x^2z^2 + x^2y^2\right)e^{xyz}
		}
		\end{solution}
	\end{minipage}\hfill%
	\begin{minipage}[t]{0.37\linewidth}
		\raggedleft \vspace{12pt}
		$\nabla^2 f$ = \ansbox[1.5in]{$(y^2z^2 + x^2z^2 + x^2y^2)e^{xyz}$}
	\end{minipage}

	\begin{minipage}[t]{0.58\linewidth}
		\part Find the Laplacian of $f = 1/(x^2 + y^2 + z^2)$. Use the coordinate
		system that makes it easiest.
		\begin{solution}[1.2in]
		In spherical coordinates $f = 1/r^2$, which depends only on $r$:
		\eq{
			\nabla^2 f = \dfrac{1}{r^2}\dfrac{\partial}{\partial r}\left(r^2 \cdot \dfrac{-2}{r^3}\right)
			= \dfrac{1}{r^2}\dfrac{\partial}{\partial r}\left(\dfrac{-2}{r}\right) = \dfrac{2}{r^4}
		}
		\end{solution}
	\end{minipage}\hfill%
	\begin{minipage}[t]{0.37\linewidth}
		\raggedleft \vspace{12pt}
		$\nabla^2 f$ = \ansbox[1.3in]{$2/r^4$}
	\end{minipage}

\end{parts}

\newpage
%%%%%%%%%%%%%%%%%%%%%%%%%%%%%%%%%%%%%%%%%
\question \textbf{The divergence theorem.} Verify the divergence theorem for
$\mathbf{F} = \hat{\boldsymbol{\rho}}\,\rho^2 + \hat{\mathbf{z}}\,z$ over the
cylinder $0 \le \rho \le a$, $0 \le z \le h$.
\begin{parts}

	\begin{minipage}[t]{0.58\linewidth}
		\part Evaluate the volume integral of $\nabla\cdot\mathbf{F}$.
		\begin{solution}[1.2in]
		From Question 2(b), $\nabla\cdot\mathbf{F} = 3\rho + 1$:
		\eq{
			\int_V \nabla\cdot\mathbf{F}\,dv &= \int_0^{h}\!\!\int_0^{2\pi}\!\!\int_0^{a}(3\rho + 1)\,\rho\,d\rho\,d\phi\,dz \\
			&= 2\pi h\left(a^3 + \dfrac{a^2}{2}\right) = 2\pi a^3h + \pi a^2h
		}
		\end{solution}
	\end{minipage}\hfill%
	\begin{minipage}[t]{0.37\linewidth}
		\raggedleft \vspace{12pt}
		\ansbox[1.5in]{$2\pi a^3h + \pi a^2h$}
	\end{minipage}

	\part Evaluate the flux of $\mathbf{F}$ out of each of the three surfaces of
	the cylinder, and show that the total equals your answer to part (a).
	\begin{solution}[2in]
	\emph{Side}, $\rho = a$: $d\mathbf{s} = \hat{\boldsymbol{\rho}}\,a\,d\phi\,dz$ and $F_\rho = a^2$:
	\eq{ \int_0^{h}\!\!\int_0^{2\pi} a^2 \cdot a\,d\phi\,dz = 2\pi a^3h }
	\emph{Top}, $z = h$: $d\mathbf{s} = \hat{\mathbf{z}}\,\rho\,d\rho\,d\phi$ and $F_z = h$:
	\eq{ \int_0^{2\pi}\!\!\int_0^{a} h\,\rho\,d\rho\,d\phi = \pi a^2h }
	\emph{Bottom}, $z = 0$: $d\mathbf{s} = -\hat{\mathbf{z}}\,\rho\,d\rho\,d\phi$ and $F_z = 0$, so the flux is 0.
	The total is $2\pi a^3h + \pi a^2h$, which matches part (a).
	\end{solution}

\end{parts}

\newpage
%%%%%%%%%%%%%%%%%%%%%%%%%%%%%%%%%%%%%%%%%
\question \textbf{Curl and Stokes' theorem.} The field
$\mathbf{F} = \hat{\boldsymbol{\phi}}\,\rho$ circulates around the $z$ axis.
\begin{parts}

	\begin{minipage}[t]{0.58\linewidth}
		\part Find $\nabla\times\mathbf{F}$.
		\begin{solution}[1in]
		Only $F_\phi$ is nonzero, and it depends only on $\rho$, so only the $z$
		component survives:
		\eq{
			(\nabla\times\mathbf{F})_z = \dfrac{1}{\rho}\dfrac{\partial\left(\rho\cdot\rho\right)}{\partial\rho} = 2
		}
		\end{solution}
	\end{minipage}\hfill%
	\begin{minipage}[t]{0.37\linewidth}
		\raggedleft \vspace{12pt}
		$\nabla\times\mathbf{F}$ = \ansbox[1.3in]{$2\hat{\mathbf{z}}$}
	\end{minipage}

	\begin{minipage}[t]{0.58\linewidth}
		\part Find the circulation of $\mathbf{F}$ counterclockwise around the circle
		$\rho = a$ in the $xy$ plane, directly and with Stokes' theorem.
		\begin{solution}[1.4in]
		Directly, $d\mathbf{l} = \hat{\boldsymbol{\phi}}\,a\,d\phi$ and $F_\phi = a$:
		\eq{ \oint_C \mathbf{F}\cdot d\mathbf{l} = \int_0^{2\pi} a \cdot a\,d\phi = 2\pi a^2 }
		With Stokes' theorem, over the disk with $d\mathbf{s} = \hat{\mathbf{z}}\,\rho\,d\rho\,d\phi$:
		\eq{ \int_S 2\,ds = 2 \cdot \pi a^2 = 2\pi a^2 }
		\end{solution}
	\end{minipage}\hfill%
	\begin{minipage}[t]{0.37\linewidth}
		\raggedleft \vspace{12pt}
		\ansbox[1.3in]{$2\pi a^2$}
	\end{minipage}

	\part Find $\nabla\cdot\mathbf{F}$. In one or two sentences, explain what the
	divergence and the curl of this field tell you about it.
	\begin{solution}[1.1in]
	$\nabla\cdot\mathbf{F} = (1/\rho)\,\partial F_\phi/\partial\phi = 0$. The field has no
	sources: every field line that enters a small volume also leaves it. Its curl
	is nonzero everywhere, so it circulates around every point, not only around
	the $z$ axis.
	\end{solution}

	\part State which theorem relates a volume integral to a closed-surface
	integral, and which relates an open-surface integral to a closed-path
	integral.
	\begin{solution}[0.8in]
	The divergence theorem relates a volume to its closed surface:
	$\int_V \nabla\cdot\mathbf{F}\,dv = \oint_S \mathbf{F}\cdot d\mathbf{s}$.
	Stokes' theorem relates an open surface to its boundary path:
	$\int_S (\nabla\times\mathbf{F})\cdot d\mathbf{s} = \oint_C \mathbf{F}\cdot d\mathbf{l}$.
	\end{solution}

\end{parts}

\vspace{1.5em}
\noindent\textbf{Documentation:} \rule[-2pt]{0.6\linewidth}{0.4pt}

\end{questions}
